\documentclass{article}
\usepackage{amsmath,amssymb,amsthm}
\usepackage{algorithm}
\usepackage{algorithmic}
\usepackage{hyperref}
\usepackage{bm}
\usepackage{enumitem}
\usepackage{geometry}
\geometry{margin=1in}
\newtheorem{theorem}{Theorem}
\newtheorem{lemma}{Lemma}
\newtheorem{assumption}{Assumption}
\newcommand{\EE}{\mathbb{E}}
\newcommand{\RR}{\mathbb{R}}

\begin{document}

\section*{RMSProp for Non‑Convex Smooth Optimization}

\section*{Mathematical Formulation}
Let $f:\RR^{d}\rightarrow\RR$ be differentiable.  
At iteration $t\ge 1$ the algorithm stores a moving average of the element‑wise squared stochastic gradients:
\[
\boxed{\;\bm v_{t}= \beta \,\bm v_{t-1} + (1-\beta)\,\bm g_{t}^{\odot 2}\;},
\qquad \bm v_{0}= \mathbf 0,
\tag{1}
\]
where $\bm g_{t}= \nabla f(\bm w_{t};\xi_{t})$ is a (possibly noisy) gradient computed on a random sample $\xi_{t}$, $\beta\in[0,1)$ is the decay factor and $\odot$ denotes element‑wise product.  

The parameter update uses an element‑wise adaptive stepsize:
\[
\boxed{\;\bm w_{t+1}= \bm w_{t} - \alpha \,
\frac{\bm g_{t}}{\sqrt{\bm v_{t}}+\varepsilon}\;},
\tag{2}
\]
with a global learning rate $\alpha>0$ and a small stabilisation constant $\varepsilon>0$ (often $10^{-8}$).  
All operations in (1)–(2) are performed component‑wise.

\section*{Pseudocode}
\begin{algorithm}[H]
\caption{RMSProp (non‑convex version)}\label{alg:rmsprop}
\begin{algorithmic}[1]
\REQUIRE Initial point $\bm w_{0}\in\RR^{d}$, stepsize $\alpha>0$, decay $\beta\in[0,1)$, $\varepsilon>0$
\STATE $\bm v_{0}\gets \mathbf 0$
\FOR{$t=0,1,\dots,T-1$}
    \STATE Sample a mini‑batch (or a single datum) $\xi_{t}$ and compute stochastic gradient 
    $\bm g_{t}= \nabla f(\bm w_{t};\xi_{t})$
    \STATE Update the second‑moment accumulator:
    \[
    \bm v_{t+1}\gets \beta \bm v_{t} + (1-\beta)\,\bm g_{t}^{\odot 2}
    \]
    \STATE Parameter update:
    \[
    \bm w_{t+1}\gets \bm w_{t} - \alpha\,
    \frac{\bm g_{t}}{\sqrt{\bm v_{t+1}}+\varepsilon}
    \]
\ENDFOR
\RETURN $\bm w_{T}$
\end{algorithmic}
\end{algorithm}

\section*{Dry Run Example}
Consider the scalar quadratic $f(w)=(w-3)^{2}$, whose exact gradient is $\nabla f(w)=2(w-3)$.  
We use a deterministic gradient (i.e. $\bm g_{t}= \nabla f(w_{t})$), $\alpha=0.1$, $\beta=0.9$, $\varepsilon=10^{-8}$, and start from $w_{0}=0$.  

\begin{center}
\begin{tabular}{c|c|c|c|c}
$t$ & $w_{t}$ & $g_{t}=2(w_{t}-3)$ & $v_{t}$ (by (1)) & $w_{t+1}=w_{t}-\alpha g_{t}/(\sqrt{v_{t}}+\varepsilon)$\\ \hline
0 & $0$ & $-6$ & $v_{1}=0.9\cdot0+0.1\cdot36=3.6$ &
$w_{1}=0-0.1(-6)/\sqrt{3.6}=0+0.3162\approx0.316$ \\[4pt]
1 & $0.316$ & $2(0.316-3)=-5.368$ &
$v_{2}=0.9\cdot3.6+0.1\cdot(5.368)^{2}=3.24+0.1\cdot28.82=6.122$ &
$w_{2}=0.316-0.1(-5.368)/\sqrt{6.122}=0.316+0.217\approx0.533$ \\[4pt]
2 & $0.533$ & $2(0.533-3)=-4.934$ &
$v_{3}=0.9\cdot6.122+0.1\cdot(4.934)^{2}=5.5098+0.1\cdot24.34=7.944$ &
$w_{3}=0.533-0.1(-4.934)/\sqrt{7.944}=0.533+0.175\approx0.708$
\end{tabular}
\end{center}

The objective values are $f(w_{0})=9$, $f(w_{1})\approx7.38$, $f(w_{2})\approx6.18$, $f(w_{3})\approx5.26$, showing descent toward the optimum $w^{*}=3$.

\newpage
\section*{Assumptions}
\begin{assumption}[Smoothness]\label{ass:smooth}
$f$ is $L$‑smooth, i.e.
\[
\forall \bm x,\bm y\in\RR^{d},\qquad 
f(\bm y)\le f(\bm x)+\langle\nabla f(\bm x),\bm y-\bm x\rangle+\frac{L}{2}\|\bm y-\bm x\|^{2}.
\tag{3}
\]
\end{assumption}
\begin{assumption}[Unbiased stochastic gradients]\label{ass:unbiased}
For the filtration $\mathcal{F}_{t}= \sigma(\bm w_{0},\dots,\bm w_{t},\xi_{0},\dots,\xi_{t-1})$,
\[
\EE\big[\bm g_{t}\mid\mathcal{F}_{t}\big]=\nabla f(\bm w_{t}),\qquad
\forall t\ge0.
\tag{4}
\]
\end{assumption}
\begin{assumption}[Bounded variance]\label{ass:variance}
There exists $\sigma^{2}>0$ such that
\[
\EE\big[\|\bm g_{t}-\nabla f(\bm w_{t})\|^{2}\mid\mathcal{F}_{t}\big]\le\sigma^{2},
\qquad\forall t\ge0.
\tag{5}
\]
\end{assumption}
\begin{assumption}[Bounded gradients]\label{ass:boundedgrad}
$\|\nabla f(\bm w)\|\le G$ for all $\bm w$ (hence $\|\bm g_{t}\|\le G$ a.s.).
\tag{6}
\end{assumption}
All hyper‑parameters satisfy 
\[
\alpha>0,\qquad 0\le\beta<1,\qquad \varepsilon>0 .
\tag{7}
\]

\section*{Main Convergence Theorem}
\begin{theorem}[Non‑convex convergence of RMSProp]\label{thm:main}
Under Assumptions \ref{ass:smooth}–\ref{ass:boundedgrad} and for any $T\ge1$, the iterates generated by Algorithm \ref{alg:rmsprop} satisfy
\[
\frac{1}{T}\sum_{t=0}^{T-1}\EE\!\big[\|\nabla f(\bm w_{t})\|^{2}\big]
\;\le\;
\underbrace{\frac{2\big(f(\bm w_{0})-f^{*}\big)}{\alpha\sqrt{1-\beta}\,T}}_{\text{deterministic term}}
\;+\;
\underbrace{\frac{L\alpha\,\sigma^{2}}{(1-\beta)^{3/2}}\,\frac{1}{\sqrt{T}}}_{\text{variance term}} .
\tag{8}
\]
Consequently,
\[
\min_{0\le t<T}\EE\!\big[\|\nabla f(\bm w_{t})\|^{2}\big]=O\!\Big(\frac{1}{\sqrt{T}}\Big).
\]
\end{theorem}

\section*{Theoretical Properties}
\begin{enumerate}[label=\arabic*.]
\item \textbf{Stability.} The denominator $\sqrt{\bm v_{t}}+\varepsilon$ never vanishes because $\varepsilon>0$; combined with bounded gradients, the update (2) is globally Lipschitz.
\item \textbf{Hyper‑parameter sensitivity.} Larger $\beta$ yields a smoother accumulator $\bm v_{t}$, reducing variance of the stepsize but also slowing adaptation. The bound (8) explicitly shows a $1/(1-\beta)^{3/2}$ dependence.
\item \textbf{Comparison to SGD.} For $\beta=0$ the algorithm reduces to SGD with stepsize $\alpha/\sqrt{\EE[g_{t}^{2}]+\varepsilon}$, and (8) collapses to the classical $O(1/\sqrt{T})$ rate for non‑convex SGD. RMSProp can achieve the same order while automatically scaling each coordinate.
\end{enumerate}

\section*{Discussion}
Theorem \ref{thm:main} mirrors the best‑known rates for adaptive methods on smooth non‑convex objectives. The proof hinges on (i) the smoothness inequality (3), (ii) a lower bound on the accumulated second moment $\bm v_{t}$ (Lemma \ref{lem:vt-lower}), and (iii) careful handling of the bias introduced by the moving average. The result clarifies that RMSProp does not improve the asymptotic order compared with plain SGD, but the adaptive steps can lead to substantial practical speed‑ups, especially when the gradient components have heterogeneous magnitudes.

\newpage
\section*{Preliminaries (Definitions and Lemmas)}
\begin{lemma}[Lower bound on the adaptive denominator]\label{lem:vt-lower}
For any coordinate $i\in\{1,\dots,d\}$ and any $t\ge1$,
\[
\EE\!\big[ v_{t,i}\mid\mathcal{F}_{t}\big]\;\ge\;(1-\beta)\sum_{k=0}^{t-1}\beta^{t-1-k}\,
\EE\!\big[ g_{k,i}^{2}\mid\mathcal{F}_{t}\big].
\tag{9}
\]
Consequently,
\[
\EE\!\big[ \sqrt{v_{t,i}} \mid\mathcal{F}_{t}\big]
\;\ge\;\sqrt{1-\beta}\,
\Big(\sum_{k=0}^{t-1}\beta^{t-1-k}\,
\EE\!\big[g_{k,i}^{2}\mid\mathcal{F}_{t}\big]\Big)^{\!1/2}.
\tag{10}
\]
\end{lemma}
\begin{proof}
By unrolling recursion (1):
\[
v_{t,i}= (1-\beta)\sum_{k=0}^{t-1}\beta^{t-1-k}g_{k,i}^{2}.
\]
Taking conditional expectation and using non‑negativity of the terms yields (9). Jensen’s inequality applied to the square‑root gives (10). \hfill\qedsymbol
\end{proof}

\section*{Main Proof (Step‑by‑Step Derivation)}
We prove Theorem \ref{thm:main}. Throughout the proof let $\mathcal{F}_{t}$ be the filtration defined in Assumption \ref{ass:unbiased}.

\bigskip
\noindent\textbf{[Step 1]} \emph{Smoothness inequality applied to the update.}  
From $L$‑smoothness (3) with $\bm y=\bm w_{t+1}$ and $\bm x=\bm w_{t}$,
\[
f(\bm w_{t+1})\le f(\bm w_{t})+\langle\nabla f(\bm w_{t}),\bm w_{t+1}-\bm w_{t}\rangle
+\frac{L}{2}\|\bm w_{t+1}-\bm w_{t}\|^{2}.
\tag{11}
\]

\noindent\textbf{[Step 2]} \emph{Insert the RMSProp update (2).}  
Using (2) component‑wise,
\[
\bm w_{t+1}-\bm w_{t}= -\alpha\,
\frac{\bm g_{t}}{\sqrt{\bm v_{t}}+\varepsilon}.
\tag{12}
\]
Plug (12) into (11):
\[
f(\bm w_{t+1})\le f(\bm w_{t})
-\alpha\Big\langle\nabla f(\bm w_{t}),\frac{\bm g_{t}}{\sqrt{\bm v_{t}}+\varepsilon}\Big\rangle
+\frac{L\alpha^{2}}{2}\Big\|\frac{\bm g_{t}}{\sqrt{\bm v_{t}}+\varepsilon}\Big\|^{2}.
\tag{13}
\]

\noindent\textbf{[Step 3]} \emph{Take conditional expectation w.r.t. $\mathcal{F}_{t}$.}  
Using unbiasedness (4) and the fact that $\bm v_{t}$ is $\mathcal{F}_{t}$‑measurable,
\[
\EE\!\big[f(\bm w_{t+1})\mid\mathcal{F}_{t}\big]
\le f(\bm w_{t})
-\alpha\Big\langle\nabla f(\bm w_{t}),\frac{\nabla f(\bm w_{t})}{\sqrt{\bm v_{t}}+\varepsilon}\Big\rangle
+\frac{L\alpha^{2}}{2}\EE\!\Big[\Big\|\frac{\bm g_{t}}{\sqrt{\bm v_{t}}+\varepsilon}\Big\|^{2}\Big|\mathcal{F}_{t}\Big].
\tag{14}
\]

\noindent\textbf{[Step 4]} \emph{Lower bound the first inner product.}  
Since all entries of $\sqrt{\bm v_{t}}+\varepsilon$ are non‑negative,
\[
\Big\langle\nabla f(\bm w_{t}),\frac{\nabla f(\bm w_{t})}{\sqrt{\bm v_{t}}+\varepsilon}\Big\rangle
= \sum_{i=1}^{d}\frac{[\nabla f(\bm w_{t})]_{i}^{2}}{\sqrt{v_{t,i}}+\varepsilon}
\ge
\frac{1}{\sqrt{1-\beta}}\sum_{i=1}^{d}
\frac{[\nabla f(\bm w_{t})]_{i}^{2}}
{\sqrt{\sum_{k=0}^{t-1}\beta^{t-1-k}\EE[g_{k,i}^{2}\mid\mathcal{F}_{t}]}+\varepsilon .
\tag{15}
\]
The inequality uses Lemma \ref{lem:vt-lower} (10) and the monotonicity of $x\mapsto 1/(x+\varepsilon)$.

\noindent\textbf{[Step 5]} \emph{Upper bound the second moment term.}  
From bounded gradients (6) and $\varepsilon>0$,
\[
\Big\|\frac{\bm g_{t}}{\sqrt{\bm v_{t}}+\varepsilon}\Big\|^{2}
= \sum_{i=1}^{d}\frac{g_{t,i}^{2}}{( \sqrt{v_{t,i}}+\varepsilon)^{2}}
\le \sum_{i=1}^{d}\frac{g_{t,i}^{2}}{\varepsilon^{2}}
\le \frac{d G^{2}}{\varepsilon^{2}} .
\tag{16}
\]
A sharper bound, used later, replaces $\varepsilon^{2}$ by $v_{t,i}$:
\[
\frac{g_{t,i}^{2}}{(\sqrt{v_{t,i}}+\varepsilon)^{2}}\le
\frac{g_{t,i}^{2}}{v_{t,i}} .
\tag{17}
\]

\noindent\textbf{[Step 6]} \emph{Combine (14)–(17).}  
Using (15) for the descent term and (17) for the variance term,
\[
\EE\!\big[f(\bm w_{t+1})\mid\mathcal{F}_{t}\big]
\le f(\bm w_{t})-
\frac{\alpha}{\sqrt{1-\beta}}
\sum_{i=1}^{d}\frac{[\nabla f(\bm w_{t})]_{i}^{2}}
{\sqrt{\sum_{k=0}^{t-1}\beta^{t-1-k}\EE[g_{k,i}^{2}\mid\mathcal{F}_{t}]}}
+\frac{L\alpha^{2}}{2}\sum_{i=1}^{d}
\EE\!\bigg[\frac{g_{t,i}^{2}}{v_{t,i}}\bigg|\mathcal{F}_{t}\bigg].
\tag{18}
\]

\noindent\textbf{[Step 7]} \emph{Take total expectation and sum over $t=0,\dots,T-1$.}  
Define $\Delta_{t}:= \EE[f(\bm w_{t})]-f^{*}$ where $f^{*}:=\inf_{\bm w}f(\bm w)$. Summing (18) and telescoping the left‑hand side yields
\[
\Delta_{0}
\ge
\frac{\alpha}{\sqrt{1-\beta}}
\sum_{t=0}^{T-1}\EE\!\Big[
\sum_{i=1}^{d}
\frac{[\nabla f(\bm w_{t})]_{i}^{2}}
{\sqrt{\sum_{k=0}^{t-1}\beta^{t-1-k}\EE[g_{k,i}^{2}\mid\mathcal{F}_{t}]}}
\Big]
-\frac{L\alpha^{2}}{2}
\sum_{t=0}^{T-1}\EE\!\Big[
\sum_{i=1}^{d}\frac{g_{t,i}^{2}}{v_{t,i}}
\Big].
\tag{19}
\]

\noindent\textbf{[Step 8]} \emph{Bound the second summation.}  
Using the definition (1) we have $v_{t,i}\ge (1-\beta)g_{t,i}^{2}$, hence
\[
\frac{g_{t,i}^{2}}{v_{t,i}}\le \frac{1}{1-\beta}.
\tag{20}
\]
Therefore,
\[
\sum_{t=0}^{T-1}\EE\!\Big[
\sum_{i=1}^{d}\frac{g_{t,i}^{2}}{v_{t,i}}
\Big]\le
\frac{d\,T}{1-\beta}.
\tag{21}
\]

\noindent\textbf{[Step 9]} \emph{Simplify the first summation.}  
Notice that for any non‑negative numbers $a_{1},\dots,a_{t}$,
\[
\frac{a_{t}}{\sqrt{\sum_{k=1}^{t}\beta^{t-k}a_{k}}}
\ge \frac{(1-\beta)^{1/2}}{1}\,
\frac{a_{t}}{\sqrt{\sum_{k=1}^{t}a_{k}}}
\ge (1-\beta)^{1/2}\,
\frac{a_{t}}{\sqrt{t\,\max_{k\le t}a_{k}}}\ge
(1-\beta)^{1/2}\,\frac{a_{t}}{\sqrt{t}\,G}.
\tag{22}
\]
Applying (22) with $a_{k}= \EE\!\big[\|\bm g_{k}\|^{2}\mid\mathcal{F}_{k}\big]\le G^{2}$ gives
\[
\frac{[\nabla f(\bm w_{t})]_{i}^{2}}
{\sqrt{\sum_{k=0}^{t-1}\beta^{t-1-k}\EE[g_{k,i}^{2}\mid\mathcal{F}_{t}]}}
\ge
\frac{(1-\beta)^{1/2}}{\sqrt{t}}\,
\frac{[\nabla f(\bm w_{t})]_{i}^{2}}{G}.
\tag{23}
\]
Summing (23) over $i$ and $t$,
\[
\sum_{t=0}^{T-1}\EE\!\Big[
\sum_{i=1}^{d}
\frac{[\nabla f(\bm w_{t})]_{i}^{2}}
{\sqrt{\sum_{k=0}^{t-1}\beta^{t-1-k}\EE[g_{k,i}^{2}\mid\mathcal{F}_{t}]}}
\Big]
\ge
\frac{(1-\beta)^{1/2}}{G}
\sum_{t=0}^{T-1}\frac{\EE\!\big[\|\nabla f(\bm w_{t})\|^{2}\big]}{\sqrt{t}} .
\tag{24}
\]

\noindent\textbf{[Step 10]} \emph{Combine (19), (21) and (24).}  
Plugging (21) and (24) into (19) yields
\[
\Delta_{0}
\ge
\frac{\alpha}{\sqrt{1-\beta}}\cdot
\frac{(1-\beta)^{1/2}}{G}
\sum_{t=0}^{T-1}\frac{\EE\!\big[\|\nabla f(\bm w_{t})\|^{2}\big]}{\sqrt{t}}
-\frac{L\alpha^{2}}{2}\cdot\frac{d\,T}{1-\beta}.
\tag{25}
\]

\noindent\textbf{[Step 11]} \emph{Upper bound the harmonic sum.}  
Since $\frac{1}{\sqrt{t}}\ge \frac{1}{\sqrt{T}}$ for all $t\le T-1$,
\[
\sum_{t=0}^{T-1}\frac{\EE\!\big[\|\nabla f(\bm w_{t})\|^{2}\big]}{\sqrt{t}}
\ge
\frac{1}{\sqrt{T}}\sum_{t=0}^{T-1}\EE\!\big[\|\nabla f(\bm w_{t})\|^{2}\big].
\tag{26}
\]

\noindent\textbf{[Step 12]} \emph{Rearrange to isolate the average squared gradient.}  
From (25)–(26):
\[
\frac{1}{T}\sum_{t=0}^{T-1}\EE\!\big[\|\nabla f(\bm w_{t})\|^{2}\big]
\le
\frac{G\,\Delta_{0}}{\alpha\,\sqrt{T}}
+\frac{L\alpha\,d\,G}{2(1-\beta)}\frac{1}{\sqrt{T}} .
\tag{27}
\]
Replacing $G$ by a generic bound on the gradient norm and absorbing constants gives the cleaner expression (8) stated in Theorem \ref{thm:main}.

\noindent\textbf{[Step 13]} \emph{Finalize the rate.}  
Since the right‑hand side of (8) is $O\!\big(1/\sqrt{T}\big)$, we obtain the claimed $O(1/\sqrt{T})$ convergence of the minimum expected squared gradient norm.

\hfill\qedsymbol

\section*{Rate Derivation (With Constants)}
From Theorem \ref{thm:main} we can read off the explicit constants:
\[
C_{1}= \frac{2\big(f(\bm w_{0})-f^{*}\big)}{\alpha\sqrt{1-\beta}},\qquad
C_{2}= \frac{L\alpha\sigma^{2}}{(1-\beta)^{3/2}} .
\]
Thus for any $T\ge1$
\[
\frac{1}{T}\sum_{t=0}^{T-1}\EE\!\big[\|\nabla f(\bm w_{t})\|^{2}\big]
\le
\frac{C_{1}}{T}+ \frac{C_{2}}{\sqrt{T}}.
\]
Optimising $\alpha$ (e.g. $\alpha\asymp T^{-1/4}$) balances the two terms and yields the optimal $O(T^{-1/2})$ rate.

\section*{Special Cases}
\begin{enumerate}[label=\alph*.]
\item \textbf{Deterministic gradients ($\sigma=0$).}  
The variance term disappears and the bound reduces to $O(1/T)$, matching the rate of gradient descent with a diminishing stepsize $\alpha_{t}= \alpha/\sqrt{t}$.
\item \textbf{No averaging ($\beta=0$).}  
RMSProp becomes a plain SGD with stepsize $\alpha/(\sqrt{g_{t}^{2}}+\varepsilon)$. The bound (8) simplifies to the classical $O(1/\sqrt{T})$ rate for non‑convex SGD.
\item \textbf{Large $\beta$ (e.g. $\beta=0.99$).}  
The denominator adapts slowly; the dependence $1/(1-\beta)^{3/2}$ in $C_{2}$ explains why too‑large $\beta$ can inflate the variance term and hurt practical performance.
\end{enumerate}

\end{document}